\documentclass[10pt,a4paper]{amsart}
\usepackage[margin=19mm]{geometry}
\usepackage{amsmath,amssymb,mathtools}
\usepackage[scr=rsfs,cal=euler]{mathalfa}
\usepackage{xfrac}
\usepackage[unicode,hidelinks,bookmarks=false]{hyperref}
\newcommand{\ie}{i.e.\ }
\newcommand{\cf}{cf.\ }
\newcommand{\resp}{respectively\ }
\newcommand{\Subsection}{\subsection}
\providecommand{\nobreakdash}{\nobreak\hbox{-}}
\setlength{\emergencystretch}{2em}
\multlinegap0pt
\usepackage{amssymb}
\usepackage[shortlabels]{enumitem}
\newtheorem{proposition}{Proposition}
\title{Operators arising from invariant measures under some class of multidimensional transformations}
\author{Oleksandr V. Maslyuchenko, Janusz Morawiec, Thomas Z\"urcher}
\date{Source: arXiv:2603.06076v1 (CC BY 4.0)}
\begin{document}
\maketitle

\noindent\emph{Source and licence.} Operators arising from invariant measures under some class of multidimensional transformations. Version arXiv:2603.06076v1 (CC BY 4.0), distributed by arXiv under the Creative Commons Attribution 4.0 International licence, http://creativecommons.org/licenses/by/4.0/, as recorded in the arXiv licence metadata for this exact version and shown on the abstract page https://arxiv.org/abs/2603.06076v1. Full licence notice: \texttt{licenses/arxiv-2603.06076-CC-BY-4.0.txt}.\par\smallskip

\noindent\textit{Editorial scope.} Counted unit: the article's proposition prop:IterationsMWoperatorsSpecial (source0.tex lines 577--591). The article's complete original proof of it is delivered with it (source0.tex lines 593--674, one \texttt{proof} environment of 82 lines). No mathematics is written, repaired or re-derived: the statement and the proof are the article's own lines, in the article's own notation, and no retained line is substituted or re-flowed.  Retained uncounted as the source-local context the proof needs: lines 444--446; lines 458--460.  Nothing was omitted from any retained span, so the package carries no omission record.  No retained line contains a citation, so the wrapper carries no bibliography and no bibliography entry.  No retained line contains a figure environment, an image-inclusion command or drawing code, so no image is delivered and the package declares no asset dependency.  Editorial formatting only, in addition: an AMS \texttt{amsart} preamble that restates the article's own declarations and macros used by the retained lines, namely the environments \texttt{proposition} on separate counters, together with the standard packages those lines need.  The retained environments print in this document as Proposition 1 in order of appearance, whereas the article prints the same items with the numbering of its own full document.  The complete native source of the article as distributed in the arXiv e-print for this exact version is bundled unchanged as \texttt{sources/195839/source0.tex}.
\begin{equation}\label{gammaiM}
\gamma^i_M(x)=\pi_M\big(\gamma^i(0),\gamma^i(x)\big),
\end{equation}

\begin{equation}\label{equ:MWoperatorDirect}
\mathbb{M}f(x)=\sum_{i\in I}\sum_{M\subseteq N}(-1)^{|M|}f\big(\gamma^i_M(x)\big).
\end{equation}

\begin{proposition}\label{prop:IterationsMWoperatorsSpecial}
Assume that for any $i\in I$, the function $\gamma^i$ is of the form
\begin{equation*}
\gamma^i(x)=\big(\gamma^i_n(x_n)\big)_{n\in N}\quad\text{for every }x=(x_n)_{n\in N}\in D,
\end{equation*}	
where $\gamma^i_n\colon D_n\to D_n$ and $D_n=\{x_n\in V:x\in D\}$ for every $n\in N$.
Then for all $p\in\mathbb{N}$ and $f\in \mathcal{F}^p(X)$ we have 
\begin{equation*}\label{equ:mMWoperatorSpecial}
\mathbb{M}^p f(x)=\sum_{i\in I^p}\square f\big(\gamma^i(0),\gamma^i(x)\big),
\end{equation*}
or equivalently, extending the  definition of $\gamma^i_M$ given in \eqref{gammaiM} to all $i \in I^p$, $x \in X$, and $M \subseteq N$, 
\begin{equation}\label{equ:mMWoperatorSpecialDirect}
\mathbb{M}^p f(x)=\sum_{i\in I^p}\sum_{M\subseteq N}(-1)^{|M|}f\big(\gamma^i_M(x)\big).
\end{equation}
\end{proposition}	

\begin{proof}
We proceed by induction on $p$ to prove \eqref{equ:mMWoperatorSpecialDirect}.
For $p=1$, \eqref{equ:mMWoperatorSpecialDirect} reduces to~\eqref{equ:MWoperatorDirect}. Fix $p\in\mathbb N$ and suppose that \eqref{equ:mMWoperatorSpecialDirect} holds for every $f\in \mathcal{F}^p(X)$.

Let us fix $f\in\mathcal{F}^{p+1}(X)$ and $x\in X$. Then $\mathbb{M}f\in\mathcal{F}^p(X)$. Putting 
\begin{equation*}
\Sigma_M^{i,j}(x)=\sum_{L\subseteq N}(-1)^{|L|}f\big(\gamma^i_M(\gamma^j_L(x))\big)
\end{equation*}
for all $i\in I$, $j\in I^p$ and $M\subseteq N$, we have
\begin{equation}\label{equ:MWspecial(p+1)}
\begin{split}
\mathbb{M}^{p+1} f(x)&=\mathbb{M}^p\big(\mathbb{M} f\big)(x)
=\sum_{j\in I^p}\sum_{L\subseteq N}(-1)^{|L|}\mathbb{M}f\big(\gamma^j_L(x)\big)\\
&=\sum_{j\in I^p}\sum_{L\subseteq N}(-1)^{|L|}
\sum_{i\in I}\sum_{M\subseteq N}(-1)^{|M|}f\big(\gamma^i_M(\gamma^j_L(x))\big)\\
&=\sum_{j\in I^p}\sum_{i\in I}\sum_{M\subseteq N}(-1)^{|M|}\Sigma_M^{i,j}(x).
\end{split}
\end{equation}

Let $M\subseteq N$ and $L=L'\cup L''$ with $L'\subseteq M$ and $L''\subseteq N\setminus M$. Then
\begin{equation}\label{equ:gammaKLij}
\begin{split}
\gamma^i_{M,n}(\gamma^j_L(x))
&=\begin{cases}
\gamma^i_n(0),&\text{if }n\in M,\\
\gamma^i_n\big(\gamma^j_{L,n}(x_n)\big),&\text{if }n\in  N\setminus M
\end{cases}\\
&=\begin{cases}
\gamma^i_n(0),&\text{if }n\in M,\\
\gamma^i_n\big(\gamma^j_n(0)\big),&\text{if }n\in L'',\\
\gamma^i_n\big(\gamma^j_n(x_n)\big),&\text{if }n\in (N\setminus M)\setminus L''.
\end{cases}
\end{split}
\end{equation}
Therefore, $\gamma^i_M\big(\gamma^j_L(x)\big)$ does not depend on~$L'$ and
$\gamma^i_M\big(\gamma^j_L(x)\big)=\gamma^i_M\big(\gamma^j_{L''}(x)\big)$.
In consequence, 
\begin{equation}\label{Sigma}
\begin{split}
\Sigma_M^{i,j}(x)&=\sum_{L'\subseteq M}(-1)^{|L'|}\sum_{L''\subseteq N\setminus M}(-1)^{|L''|}
f\big(\gamma^i_M(\gamma^j_{L'\cup L''}(x))\big)\\
&=\sum_{L'\subseteq M}(-1)^{|L'|}\sum_{L''\subseteq N\setminus M}
(-1)^{|L''|}f\big(\gamma^i_M(\gamma^j_{ L''}(x))\big).	
\end{split}		
\end{equation}

Observe now that
\begin{equation}\label{sumL'K}
\sum_{L'\subseteq M}(-1)^{|L'|}=
\begin{cases}
0,&\text{if }M\neq\emptyset,\\
1,&\text{if }M=\emptyset. 
\end{cases}
\end{equation}

For any $i\in I$ and $j=(j_k)_{k=1}^p\in I^p$ we put $i\smallsmile j=(i,j_1,\dots,j_p)\in I^{p+1}$.
Applying~\eqref{equ:gammaKLij} we obtain 
\begin{equation*}
\gamma^i_{\emptyset,n}(\gamma^j_L(x))
=\begin{cases}
\gamma^i_n\big(\gamma^j_n(0)\big),&\text{if }n\in L,\\
\gamma^i_n\big(\gamma^j_n(x_n)\big),&\text{if }n\in N\setminus L
\end{cases}
=\gamma^{i\smallsmile j}_{L,n}(x_n)
\end{equation*}
for every $n\in N$, and hence $\gamma^i_{\emptyset}(\gamma^j_L(x))=\gamma^{i\smallsmile j}_{L}(x)$.
This together with \eqref{sumL'K} lets us continue the computation started in \eqref{Sigma} to get
\begin{equation*}
\Sigma^{i,j}_M(x)=
\begin{cases}
0,&\text{if }M\ne\emptyset,\\ 			
\sum_{L\subseteq N}(-1)^{|L|}f\big(\gamma^{i\smallsmile j}_L(x)\big)&\text{if }M=\emptyset.
\end{cases}
\end{equation*}
Finally, returning to \eqref{equ:MWspecial(p+1)} and substituting $i\smallsmile j$ by $\ell$, we conclude that
\begin{align*}
\mathbb{M}^{p+1} f(x)&=\sum_{i\in I,j\in  I^{p}}\Sigma_\emptyset^{i,j}(x)
=\sum_{i\in I, j\in  I^{p}}\sum_{L\subseteq N}(-1)^{|L|}f\big(\gamma^{i\smallsmile j}_L(x)\big)\\
&=\sum_{\ell\in I^{p+1}}\sum_{L\subseteq N}(-1)^{|L|}f\big(\gamma^{\ell}_L(x)\big),
\end{align*}
and the proof is complete.
\end{proof}

\end{document}
